% Options for packages loaded elsewhere
\PassOptionsToPackage{unicode}{hyperref}
\PassOptionsToPackage{hyphens}{url}
\PassOptionsToPackage{dvipsnames,svgnames,x11names}{xcolor}
%
\documentclass[
  10pt,
]{article}
\usepackage{amsmath,amssymb}
\usepackage{iftex}
\ifPDFTeX
  \usepackage[T1]{fontenc}
  \usepackage[utf8]{inputenc}
  \usepackage{textcomp} % provide euro and other symbols
\else % if luatex or xetex
  \usepackage{unicode-math} % this also loads fontspec
  \defaultfontfeatures{Scale=MatchLowercase}
  \defaultfontfeatures[\rmfamily]{Ligatures=TeX,Scale=1}
\fi
\ifPDFTeX\else
  % xetex/luatex font selection
  \setmainfont[]{DejaVu Serif}
  \setmonofont[]{DejaVu Sans Mono}
\fi
% Use upquote if available, for straight quotes in verbatim environments
\IfFileExists{upquote.sty}{\usepackage{upquote}}{}
\IfFileExists{microtype.sty}{% use microtype if available
  \usepackage[]{microtype}
  \UseMicrotypeSet[protrusion]{basicmath} % disable protrusion for tt fonts
}{}
\makeatletter
\@ifundefined{KOMAClassName}{% if non-KOMA class
  \IfFileExists{parskip.sty}{%
    \usepackage{parskip}
  }{% else
    \setlength{\parindent}{0pt}
    \setlength{\parskip}{6pt plus 2pt minus 1pt}}
}{% if KOMA class
  \KOMAoptions{parskip=half}}
\makeatother
\usepackage{xcolor}
\usepackage[margin=2.2cm]{geometry}
\usepackage{color}
\usepackage{fancyvrb}
\newcommand{\VerbBar}{|}
\newcommand{\VERB}{\Verb[commandchars=\\\{\}]}
\DefineVerbatimEnvironment{Highlighting}{Verbatim}{commandchars=\\\{\}}
% Add ',fontsize=\small' for more characters per line
\usepackage{framed}
\definecolor{shadecolor}{RGB}{248,248,248}
\newenvironment{Shaded}{\begin{snugshade}}{\end{snugshade}}
\newcommand{\AlertTok}[1]{\textcolor[rgb]{0.94,0.16,0.16}{#1}}
\newcommand{\AnnotationTok}[1]{\textcolor[rgb]{0.56,0.35,0.01}{\textbf{\textit{#1}}}}
\newcommand{\AttributeTok}[1]{\textcolor[rgb]{0.13,0.29,0.53}{#1}}
\newcommand{\BaseNTok}[1]{\textcolor[rgb]{0.00,0.00,0.81}{#1}}
\newcommand{\BuiltInTok}[1]{#1}
\newcommand{\CharTok}[1]{\textcolor[rgb]{0.31,0.60,0.02}{#1}}
\newcommand{\CommentTok}[1]{\textcolor[rgb]{0.56,0.35,0.01}{\textit{#1}}}
\newcommand{\CommentVarTok}[1]{\textcolor[rgb]{0.56,0.35,0.01}{\textbf{\textit{#1}}}}
\newcommand{\ConstantTok}[1]{\textcolor[rgb]{0.56,0.35,0.01}{#1}}
\newcommand{\ControlFlowTok}[1]{\textcolor[rgb]{0.13,0.29,0.53}{\textbf{#1}}}
\newcommand{\DataTypeTok}[1]{\textcolor[rgb]{0.13,0.29,0.53}{#1}}
\newcommand{\DecValTok}[1]{\textcolor[rgb]{0.00,0.00,0.81}{#1}}
\newcommand{\DocumentationTok}[1]{\textcolor[rgb]{0.56,0.35,0.01}{\textbf{\textit{#1}}}}
\newcommand{\ErrorTok}[1]{\textcolor[rgb]{0.64,0.00,0.00}{\textbf{#1}}}
\newcommand{\ExtensionTok}[1]{#1}
\newcommand{\FloatTok}[1]{\textcolor[rgb]{0.00,0.00,0.81}{#1}}
\newcommand{\FunctionTok}[1]{\textcolor[rgb]{0.13,0.29,0.53}{\textbf{#1}}}
\newcommand{\ImportTok}[1]{#1}
\newcommand{\InformationTok}[1]{\textcolor[rgb]{0.56,0.35,0.01}{\textbf{\textit{#1}}}}
\newcommand{\KeywordTok}[1]{\textcolor[rgb]{0.13,0.29,0.53}{\textbf{#1}}}
\newcommand{\NormalTok}[1]{#1}
\newcommand{\OperatorTok}[1]{\textcolor[rgb]{0.81,0.36,0.00}{\textbf{#1}}}
\newcommand{\OtherTok}[1]{\textcolor[rgb]{0.56,0.35,0.01}{#1}}
\newcommand{\PreprocessorTok}[1]{\textcolor[rgb]{0.56,0.35,0.01}{\textit{#1}}}
\newcommand{\RegionMarkerTok}[1]{#1}
\newcommand{\SpecialCharTok}[1]{\textcolor[rgb]{0.81,0.36,0.00}{\textbf{#1}}}
\newcommand{\SpecialStringTok}[1]{\textcolor[rgb]{0.31,0.60,0.02}{#1}}
\newcommand{\StringTok}[1]{\textcolor[rgb]{0.31,0.60,0.02}{#1}}
\newcommand{\VariableTok}[1]{\textcolor[rgb]{0.00,0.00,0.00}{#1}}
\newcommand{\VerbatimStringTok}[1]{\textcolor[rgb]{0.31,0.60,0.02}{#1}}
\newcommand{\WarningTok}[1]{\textcolor[rgb]{0.56,0.35,0.01}{\textbf{\textit{#1}}}}
\usepackage{longtable,booktabs,array}
\usepackage{calc} % for calculating minipage widths
% Correct order of tables after \paragraph or \subparagraph
\usepackage{etoolbox}
\makeatletter
\patchcmd\longtable{\par}{\if@noskipsec\mbox{}\fi\par}{}{}
\makeatother
% Allow footnotes in longtable head/foot
\IfFileExists{footnotehyper.sty}{\usepackage{footnotehyper}}{\usepackage{footnote}}
\makesavenoteenv{longtable}
\setlength{\emergencystretch}{3em} % prevent overfull lines
\providecommand{\tightlist}{%
  \setlength{\itemsep}{0pt}\setlength{\parskip}{0pt}}
\setcounter{secnumdepth}{-\maxdimen} % remove section numbering
\ifLuaTeX
  \usepackage{selnolig}  % disable illegal ligatures
\fi
\IfFileExists{bookmark.sty}{\usepackage{bookmark}}{\usepackage{hyperref}}
\IfFileExists{xurl.sty}{\usepackage{xurl}}{} % add URL line breaks if available
\urlstyle{same}
\hypersetup{
  pdftitle={Personal Website Manual},
  pdfauthor={Zhuang Ge},
  colorlinks=true,
  linkcolor={blue},
  filecolor={Maroon},
  citecolor={Blue},
  urlcolor={blue},
  pdfcreator={LaTeX via pandoc}}

\title{Personal Website Manual}
\usepackage{etoolbox}
\makeatletter
\providecommand{\subtitle}[1]{% add subtitle to \maketitle
  \apptocmd{\@title}{\par {\large #1 \par}}{}{}
}
\makeatother
\subtitle{Install, edit, update and deploy
https://gezhuang0717.github.io}
\author{Zhuang Ge}
\date{October 2026}

\begin{document}
\maketitle

{
\hypersetup{linkcolor=}
\setcounter{tocdepth}{3}
\tableofcontents
}
\hypertarget{what-you-have}{%
\section{1. What you have}\label{what-you-have}}

\begin{longtable}[]{@{}
  >{\raggedright\arraybackslash}p{(\columnwidth - 2\tabcolsep) * \real{0.5000}}
  >{\raggedright\arraybackslash}p{(\columnwidth - 2\tabcolsep) * \real{0.5000}}@{}}
\toprule\noalign{}
\begin{minipage}[b]{\linewidth}\raggedright
Item
\end{minipage} & \begin{minipage}[b]{\linewidth}\raggedright
Value
\end{minipage} \\
\midrule\noalign{}
\endhead
\bottomrule\noalign{}
\endlastfoot
Live site & \url{https://gezhuang0717.github.io} \\
Source repository &
\url{https://github.com/gezhuang0717/gezhuang0717.github.io} \\
Site generator & Hugo extended (version pinned to 0.167.0 in the deploy
workflow) \\
Themes & Congo (live), Blowfish, PaperMod, all installed as git
submodules \\
Deployment & GitHub Actions workflow
\texttt{.github/workflows/hugo.yml}, runs on every push to
\texttt{main} \\
\end{longtable}

How it works: you edit plain Markdown text files, push them to GitHub,
and GitHub Actions builds the HTML with Hugo and publishes it to GitHub
Pages about one minute later. Nothing needs to run on your computer
except when you want to preview.

You can edit in two ways:

\begin{itemize}
\tightlist
\item
  \textbf{In the browser only} (no installation): edit files directly on
  github.com. Section 4.
\item
  \textbf{On your Mac} (recommended for larger changes): clone, preview
  locally, push. Sections 3 and 5.
\end{itemize}

\hypertarget{repository-layout}{%
\subsection{1.1 Repository layout}\label{repository-layout}}

\begin{verbatim}
gezhuang0717.github.io/
├── content/                 ← YOUR TEXT (shared by all themes)
│   ├── _index.md            homepage bio and contact
│   ├── research.md          research themes + first-author papers
│   ├── publications.md      all publications
│   ├── projects.md
│   ├── talks.md
│   ├── mentorship.md
│   ├── cv.md
│   └── posts/               blog posts, one file each
├── assets/img/zhuang-ge.png profile picture (Congo, Blowfish)
├── static/                  files copied as-is to the site
│   └── img/zhuang-ge.png    profile picture (PaperMod)
├── sites/                   ← LOOK AND MENUS, one folder per theme
│   ├── congo/_default/      hugo.toml, params.toml, languages.en.toml, menus.en.toml
│   ├── blowfish/_default/   same files
│   └── papermod/_default/   hugo.toml (everything in one file)
├── themes/                  theme code (git submodules, do not edit)
├── .github/workflows/hugo.yml  build-and-deploy workflow
├── docs/                    this manual and short guides
├── setup.sh                 downloads the themes after cloning
└── README.md
\end{verbatim}

Rule of thumb: \textbf{text lives in \texttt{content/}, appearance lives
in \texttt{sites/\textless{}theme\textgreater{}/}}. You never need to
touch \texttt{themes/}.

\hypertarget{one-time-setup-on-your-mac}{%
\section{2. One-time setup on your
Mac}\label{one-time-setup-on-your-mac}}

\hypertarget{install-the-tools}{%
\subsection{2.1 Install the tools}\label{install-the-tools}}

Open \textbf{Terminal} (Applications → Utilities → Terminal).

\begin{enumerate}
\def\labelenumi{\arabic{enumi}.}
\item
  Install Homebrew if \texttt{brew\ -\/-version} reports ``command not
  found''. Paste the install command shown on \url{https://brew.sh} and
  follow the prompts.
\item
  Install Git and Hugo:

\begin{Shaded}
\begin{Highlighting}[]
\ExtensionTok{brew}\NormalTok{ install git hugo}
\end{Highlighting}
\end{Shaded}
\item
  Check the versions:

\begin{Shaded}
\begin{Highlighting}[]
\FunctionTok{git} \AttributeTok{{-}{-}version}
\ExtensionTok{hugo}\NormalTok{ version}
\end{Highlighting}
\end{Shaded}

  The Hugo line must contain \texttt{+extended} and a version of at
  least \textbf{0.163.0}.
\end{enumerate}

On Windows use \texttt{winget\ install\ Git.Git\ Hugo.Hugo.Extended}; on
Linux install the \texttt{hugo\_extended\_*.deb} from
\url{https://github.com/gohugoio/hugo/releases}.

\hypertarget{connect-git-to-github}{%
\subsection{2.2 Connect Git to GitHub}\label{connect-git-to-github}}

Tell Git who you are (once per computer):

\begin{Shaded}
\begin{Highlighting}[]
\FunctionTok{git}\NormalTok{ config }\AttributeTok{{-}{-}global}\NormalTok{ user.name  }\StringTok{"Zhuang Ge"}
\FunctionTok{git}\NormalTok{ config }\AttributeTok{{-}{-}global}\NormalTok{ user.email }\StringTok{"zhuang.z.ge@jyu.fi"}
\end{Highlighting}
\end{Shaded}

GitHub no longer accepts your account password for \texttt{git\ push}.
The simplest option is the GitHub CLI, which signs you in through the
browser:

\begin{Shaded}
\begin{Highlighting}[]
\ExtensionTok{brew}\NormalTok{ install gh}
\ExtensionTok{gh}\NormalTok{ auth login        }\CommentTok{\# choose GitHub.com → HTTPS → Login with a web browser}
\end{Highlighting}
\end{Shaded}

\hypertarget{download-the-site}{%
\subsection{2.3 Download the site}\label{download-the-site}}

\begin{Shaded}
\begin{Highlighting}[]
\BuiltInTok{cd}\NormalTok{ \textasciitilde{}/Documents}
\FunctionTok{git}\NormalTok{ clone }\AttributeTok{{-}{-}recursive}\NormalTok{ https://github.com/gezhuang0717/gezhuang0717.github.io}
\BuiltInTok{cd}\NormalTok{ gezhuang0717.github.io}
\end{Highlighting}
\end{Shaded}

\texttt{-\/-recursive} also downloads the three themes. If you forgot
it, run \texttt{bash\ setup.sh}.

\hypertarget{the-daily-workflow-on-your-mac}{%
\section{3. The daily workflow on your
Mac}\label{the-daily-workflow-on-your-mac}}

Every change follows the same four steps.

\textbf{Step 1 --- get the latest version} (important if you also edit
on github.com):

\begin{Shaded}
\begin{Highlighting}[]
\BuiltInTok{cd}\NormalTok{ \textasciitilde{}/Documents/gezhuang0717.github.io}
\FunctionTok{git}\NormalTok{ pull}
\end{Highlighting}
\end{Shaded}

\textbf{Step 2 --- start the live preview:}

\begin{Shaded}
\begin{Highlighting}[]
\ExtensionTok{hugo}\NormalTok{ server }\AttributeTok{{-}{-}configDir}\NormalTok{ sites/congo }\AttributeTok{{-}{-}themesDir}\NormalTok{ themes}
\end{Highlighting}
\end{Shaded}

Open \url{http://localhost:1313}. Leave this Terminal window running;
the page reloads every time you save a file.

\textbf{Step 3 --- edit} files in \texttt{content/} (or
\texttt{sites/congo/\_default/}) with any plain-text editor.
Recommended: Visual Studio Code (\url{https://code.visualstudio.com}),
then \texttt{code\ .} opens the whole folder. Save and check the
preview.

\textbf{Step 4 --- publish.} Open a second Terminal window (the first is
busy with the preview):

\begin{Shaded}
\begin{Highlighting}[]
\BuiltInTok{cd}\NormalTok{ \textasciitilde{}/Documents/gezhuang0717.github.io}
\FunctionTok{git}\NormalTok{ add .}
\FunctionTok{git}\NormalTok{ commit }\AttributeTok{{-}m} \StringTok{"Add new PRC paper"}      \CommentTok{\# short description of the change}
\FunctionTok{git}\NormalTok{ push}
\end{Highlighting}
\end{Shaded}

Then watch \textbf{Actions} on GitHub
(\url{https://github.com/gezhuang0717/gezhuang0717.github.io/actions}).
When \emph{Deploy Hugo site} shows a green tick (about 1 minute), the
live site is updated. Reload it with Cmd+Shift+R to bypass the browser
cache.

Stop the preview with \textbf{Ctrl+C}.

\hypertarget{editing-without-installing-anything}{%
\section{4. Editing without installing
anything}\label{editing-without-installing-anything}}

For a quick fix (a typo, one new paper):

\begin{enumerate}
\def\labelenumi{\arabic{enumi}.}
\tightlist
\item
  Open \url{https://github.com/gezhuang0717/gezhuang0717.github.io}.
\item
  Click into \texttt{content/}, then the file,
  e.g.~\texttt{publications.md}.
\item
  Click the \textbf{pencil icon} (Edit this file).
\item
  Make the change. Use the \textbf{Preview} tab to check the formatting.
\item
  Click \textbf{Commit changes\ldots{}}, write a short message, keep
  ``Commit directly to the main branch'', and click \textbf{Commit
  changes}.
\item
  The site rebuilds automatically; check the \textbf{Actions} tab for
  the green tick.
\end{enumerate}

To upload a file (photo, CV PDF): open the folder on github.com, click
\textbf{Add file → Upload files}, drag the file in, and commit.

If you later work on your Mac again, run \texttt{git\ pull} first so you
have these web edits.

\hypertarget{how-to-change-things}{%
\section{5. How to change things}\label{how-to-change-things}}

\hypertarget{markdown-in-two-minutes}{%
\subsection{5.1 Markdown in two minutes}\label{markdown-in-two-minutes}}

Each content file starts with a header between two \texttt{-\/-\/-}
lines. Keep it; change only \texttt{title} or \texttt{description} if
needed.

\begin{Shaded}
\begin{Highlighting}[]
\CommentTok{{-}{-}{-}}
\AnnotationTok{title:}\CommentTok{ "Research"}
\AnnotationTok{description:}\CommentTok{ "One{-}line summary shown in search results."}
\AnnotationTok{showDate:}\CommentTok{ false}
\CommentTok{{-}{-}{-}}

\FunctionTok{\#\# Section heading}
\FunctionTok{\#\#\# Sub{-}heading}
\NormalTok{Normal paragraph with **bold**, *italic* and a }\CommentTok{[}\OtherTok{link}\CommentTok{](https://doi.org/...)}\NormalTok{.}

\SpecialStringTok{{-} }\NormalTok{bullet item}
\SpecialStringTok{{-} }\NormalTok{another item}

\SpecialStringTok{1. }\NormalTok{numbered item}
\SpecialStringTok{2. }\NormalTok{next item}
\end{Highlighting}
\end{Shaded}

Isotopes: type the Unicode superscripts directly (¹⁰⁰Sn, ⁹⁷Tc, β⁻) or
use HTML:
\texttt{\textless{}sup\textgreater{}100\textless{}/sup\textgreater{}Sn}.

\hypertarget{where-each-thing-lives}{%
\subsection{5.2 Where each thing lives}\label{where-each-thing-lives}}

\begin{longtable}[]{@{}
  >{\raggedright\arraybackslash}p{(\columnwidth - 2\tabcolsep) * \real{0.5000}}
  >{\raggedright\arraybackslash}p{(\columnwidth - 2\tabcolsep) * \real{0.5000}}@{}}
\toprule\noalign{}
\begin{minipage}[b]{\linewidth}\raggedright
You want to change
\end{minipage} & \begin{minipage}[b]{\linewidth}\raggedright
Edit this file
\end{minipage} \\
\midrule\noalign{}
\endhead
\bottomrule\noalign{}
\endlastfoot
Homepage text, contact line & \texttt{content/\_index.md} \\
Research themes & \texttt{content/research.md} \\
Publications & \texttt{content/publications.md} \\
Projects / Talks / Mentorship / CV & \texttt{content/projects.md},
\texttt{talks.md}, \texttt{mentorship.md}, \texttt{cv.md} \\
Blog & \texttt{content/posts/*.md} \\
Name, headline under the photo, social icons &
\texttt{sites/congo/\_default/languages.en.toml} \\
Menu (order, names, new items) &
\texttt{sites/congo/\_default/menus.en.toml} \\
Colours, homepage layout, header style &
\texttt{sites/congo/\_default/params.toml} \\
Site title, base settings & \texttt{sites/congo/\_default/hugo.toml} \\
Profile photo & \texttt{assets/img/zhuang-ge.png} and
\texttt{static/img/zhuang-ge.png} \\
Which theme is live & \texttt{THEME:} in
\texttt{.github/workflows/hugo.yml} \\
\end{longtable}

If you switch the live theme to Blowfish or PaperMod, edit the matching
folder under \texttt{sites/} instead of \texttt{sites/congo/}.

\hypertarget{add-a-new-publication}{%
\subsection{5.3 Add a new publication}\label{add-a-new-publication}}

Open \texttt{content/publications.md}. First-author papers are a
numbered list, newest first. Insert the new line at the top and renumber
(Markdown renumbers automatically if you simply start every line with
\texttt{1.}):

\begin{Shaded}
\begin{Highlighting}[]
\SpecialStringTok{1. }\NormalTok{**Z. Ge** *et al.*, Title of the paper, [Phys. Rev. C **115**, 012345 (2027)](https://doi.org/10.1103/xxxx) · }\CommentTok{[}\OtherTok{arXiv:2701.01234}\CommentTok{](https://arxiv.org/abs/2701.01234)}
\end{Highlighting}
\end{Shaded}

For a co-authored paper add a bullet under the right year in
\emph{Selected co-authored papers}:

\begin{Shaded}
\begin{Highlighting}[]
\SpecialStringTok{{-} }\NormalTok{A. Author, **Z. Ge** *et al.*, Title, [Journal **Vol**, Page (Year)](https://doi.org/...)}
\end{Highlighting}
\end{Shaded}

If it belongs to a research theme, copy the same line under that heading
in \texttt{content/research.md}.

Tip: INSPIRE-HEP gives you the reference data. Search
\texttt{fa\ Ge,\ Zhuang} on \url{https://inspirehep.net} for
first-author papers.

\hypertarget{add-a-talk}{%
\subsection{5.4 Add a talk}\label{add-a-talk}}

In \texttt{content/talks.md}, add at the top of \emph{Invited talks} or
\emph{Contributed talks}:

\begin{Shaded}
\begin{Highlighting}[]
\SpecialStringTok{{-} }\NormalTok{**2027** — *Talk title*. Conference name, dates, City, Country (}\CommentTok{[}\OtherTok{slides}\CommentTok{](https://...)}\NormalTok{)}
\end{Highlighting}
\end{Shaded}

\hypertarget{replace-the-profile-picture}{%
\subsection{5.5 Replace the profile
picture}\label{replace-the-profile-picture}}

\begin{enumerate}
\def\labelenumi{\arabic{enumi}.}
\tightlist
\item
  Prepare a square photo, at least 400 × 400 pixels, saved as PNG named
  \texttt{zhuang-ge.png}.
\item
  Replace both \texttt{assets/img/zhuang-ge.png} and
  \texttt{static/img/zhuang-ge.png}.
\item
  To use a JPG instead, name it \texttt{zhuang-ge.jpg} and change
  \texttt{image\ =\ "img/zhuang-ge.png"} in
  \texttt{sites/congo/\_default/languages.en.toml} accordingly.
\end{enumerate}

\hypertarget{put-your-cv-pdf-online}{%
\subsection{5.6 Put your CV PDF online}\label{put-your-cv-pdf-online}}

\begin{enumerate}
\def\labelenumi{\arabic{enumi}.}
\item
  Copy the PDF to \texttt{static/cv.pdf}.
\item
  In \texttt{sites/congo/\_default/menus.en.toml}, replace the CV block
  with:

\begin{Shaded}
\begin{Highlighting}[]
\KeywordTok{[[main]]}
  \DataTypeTok{name} \OperatorTok{=} \StringTok{"CV"}
  \DataTypeTok{url} \OperatorTok{=} \StringTok{"/cv.pdf"}
  \DataTypeTok{weight} \OperatorTok{=} \DecValTok{70}
\end{Highlighting}
\end{Shaded}
\end{enumerate}

The PDF is then at \url{https://gezhuang0717.github.io/cv.pdf}. To keep
the HTML CV page as well, use \texttt{name\ =\ "CV\ (PDF)"} and a
different weight instead of replacing the block.

\hypertarget{write-a-blog-post}{%
\subsection{5.7 Write a blog post}\label{write-a-blog-post}}

\begin{Shaded}
\begin{Highlighting}[]
\ExtensionTok{hugo}\NormalTok{ new content posts/low{-}q{-}decays.md }\AttributeTok{{-}{-}configDir}\NormalTok{ sites/congo }\AttributeTok{{-}{-}themesDir}\NormalTok{ themes}
\end{Highlighting}
\end{Shaded}

This creates \texttt{content/posts/low-q-decays.md} with a header. Write
the text below it and set \texttt{draft\ =\ false} (or delete the draft
line), otherwise the post is hidden on the live site. Images: put them
next to the post by turning it into a folder
(\texttt{content/posts/low-q-decays/index.md} plus \texttt{figure.png})
and reference them as \texttt{!{[}Caption{]}(figure.png)}.

\hypertarget{add-a-new-page-and-menu-item}{%
\subsection{5.8 Add a new page and menu
item}\label{add-a-new-page-and-menu-item}}

\begin{enumerate}
\def\labelenumi{\arabic{enumi}.}
\item
  Create \texttt{content/teaching.md}:

\begin{Shaded}
\begin{Highlighting}[]
\CommentTok{{-}{-}{-}}
\AnnotationTok{title:}\CommentTok{ "Teaching"}
\AnnotationTok{showDate:}\CommentTok{ false}
\CommentTok{{-}{-}{-}}
\NormalTok{Text here.}
\end{Highlighting}
\end{Shaded}
\item
  Add to \texttt{sites/congo/\_default/menus.en.toml}:

\begin{Shaded}
\begin{Highlighting}[]
\KeywordTok{[[main]]}
  \DataTypeTok{name} \OperatorTok{=} \StringTok{"Teaching"}
  \DataTypeTok{pageRef} \OperatorTok{=} \StringTok{"teaching"}
  \DataTypeTok{weight} \OperatorTok{=} \DecValTok{45}
\end{Highlighting}
\end{Shaded}
\end{enumerate}

\texttt{weight} sets the order (smaller is further left). Existing
weights are 10, 20, \ldots{} 70.

\hypertarget{social-links-under-your-name}{%
\subsection{5.9 Social links under your
name}\label{social-links-under-your-name}}

In \texttt{sites/congo/\_default/languages.en.toml}:

\begin{Shaded}
\begin{Highlighting}[]
  \DataTypeTok{links} \OperatorTok{=} \OperatorTok{[}
    \OperatorTok{\{ }\DataTypeTok{email}\OperatorTok{ =} \StringTok{"mailto:zhuang.z.ge@jyu.fi"}\OperatorTok{ \},}
    \OperatorTok{\{ }\DataTypeTok{orcid}\OperatorTok{ =} \StringTok{"https://orcid.org/0000{-}0001{-}8586{-}6134"}\OperatorTok{ \},}
    \OperatorTok{\{ }\DataTypeTok{google{-}scholar}\OperatorTok{ =} \StringTok{"https://scholar.google.com/citations?user=YOUR\_ID"}\OperatorTok{ \},}
    \OperatorTok{\{ }\DataTypeTok{researchgate}\OperatorTok{ =} \StringTok{"https://www.researchgate.net/profile/Zhuang{-}Ge"}\OperatorTok{ \},}
    \OperatorTok{\{ }\DataTypeTok{linkedin}\OperatorTok{ =} \StringTok{"https://www.linkedin.com/in/zhuang{-}ge{-}847898202"}\OperatorTok{ \},}
    \OperatorTok{\{ }\DataTypeTok{github}\OperatorTok{ =} \StringTok{"https://github.com/gezhuang0717"}\OperatorTok{ \},}
    \OperatorTok{\{ }\DataTypeTok{link}\OperatorTok{ =} \StringTok{"https://www.jyu.fi/en/people/zhuang{-}ge"}\OperatorTok{ \},}
  \OperatorTok{]}
\end{Highlighting}
\end{Shaded}

Order in the list is the order of the icons. Other icon names include
\texttt{x-twitter}, \texttt{bluesky}, \texttt{mastodon},
\texttt{youtube}, \texttt{weibo}.

\hypertarget{colours-and-layout}{%
\subsection{5.10 Colours and layout}\label{colours-and-layout}}

In \texttt{sites/congo/\_default/params.toml}:

\begin{itemize}
\tightlist
\item
  \texttt{colorScheme\ =\ "sapphire"}: alternatives \texttt{congo},
  \texttt{avocado}, \texttt{cherry}, \texttt{fire}, \texttt{ocean},
  \texttt{slate}.
\item
  \texttt{defaultAppearance\ =\ "light"} or \texttt{"dark"};
  \texttt{autoSwitchAppearance\ =\ true} follows the visitor's system
  setting.
\item
  \texttt{{[}homepage{]}\ layout\ =\ "profile"} (photo + bio) or
  \texttt{"page"} (plain text). \texttt{showRecent\ =\ true} lists
  recent blog posts under the bio.
\end{itemize}

Blowfish has more schemes (\texttt{ocean}, \texttt{forest},
\texttt{fire}, \texttt{neon}, \texttt{noir}, \texttt{autumn},
\texttt{princess}, \texttt{terminal}, \texttt{github}, \texttt{marvel},
\texttt{one-light}, \texttt{slate}, \ldots) and homepage layouts
(\texttt{profile}, \texttt{page}, \texttt{hero}, \texttt{card},
\texttt{background}, \texttt{landing}).

\hypertarget{switch-the-live-theme}{%
\subsection{5.11 Switch the live theme}\label{switch-the-live-theme}}

\begin{enumerate}
\def\labelenumi{\arabic{enumi}.}
\item
  Preview the alternative first:

\begin{Shaded}
\begin{Highlighting}[]
\ExtensionTok{hugo}\NormalTok{ server }\AttributeTok{{-}{-}configDir}\NormalTok{ sites/blowfish }\AttributeTok{{-}{-}themesDir}\NormalTok{ themes}
\ExtensionTok{hugo}\NormalTok{ server }\AttributeTok{{-}{-}configDir}\NormalTok{ sites/papermod }\AttributeTok{{-}{-}themesDir}\NormalTok{ themes}
\end{Highlighting}
\end{Shaded}
\item
  In \texttt{.github/workflows/hugo.yml} change \texttt{THEME:\ congo}
  to \texttt{blowfish} or \texttt{papermod}.
\item
  Commit and push.
\end{enumerate}

All content files stay the same; only the look changes.

\hypertarget{custom-domain-optional}{%
\subsection{5.12 Custom domain
(optional)}\label{custom-domain-optional}}

\begin{enumerate}
\def\labelenumi{\arabic{enumi}.}
\tightlist
\item
  Buy a domain and create a DNS record: \texttt{CNAME} for
  \texttt{www.yourdomain.org} → \texttt{gezhuang0717.github.io}.
\item
  Create \texttt{static/CNAME} with the single line
  \texttt{www.yourdomain.org}, commit and push.
\item
  GitHub → repository \textbf{Settings → Pages → Custom domain}: enter
  the domain and tick \textbf{Enforce HTTPS}.
\end{enumerate}

\hypertarget{keeping-the-software-up-to-date}{%
\section{6. Keeping the software up to
date}\label{keeping-the-software-up-to-date}}

\hypertarget{update-the-themes}{%
\subsection{6.1 Update the themes}\label{update-the-themes}}

\begin{Shaded}
\begin{Highlighting}[]
\FunctionTok{git}\NormalTok{ submodule update }\AttributeTok{{-}{-}remote} \AttributeTok{{-}{-}merge}
\ExtensionTok{hugo}\NormalTok{ server }\AttributeTok{{-}{-}configDir}\NormalTok{ sites/congo }\AttributeTok{{-}{-}themesDir}\NormalTok{ themes   }\CommentTok{\# check nothing broke}
\FunctionTok{git}\NormalTok{ add themes}
\FunctionTok{git}\NormalTok{ commit }\AttributeTok{{-}m} \StringTok{"Update themes"}
\FunctionTok{git}\NormalTok{ push}
\end{Highlighting}
\end{Shaded}

Theme updates occasionally rename a setting; if the preview shows an
error, read the message, check the theme's documentation, or roll back
with \texttt{git\ checkout\ themes}.

\hypertarget{update-hugo}{%
\subsection{6.2 Update Hugo}\label{update-hugo}}

\begin{itemize}
\tightlist
\item
  On your Mac: \texttt{brew\ upgrade\ hugo}.
\item
  On GitHub: change \texttt{HUGO\_VERSION:\ 0.167.0} in
  \texttt{.github/workflows/hugo.yml} to the new version number from
  \url{https://github.com/gohugoio/hugo/releases}.
\end{itemize}

Keep the two roughly in step so the preview matches the live site.

\hypertarget{undoing-a-mistake}{%
\section{7. Undoing a mistake}\label{undoing-a-mistake}}

\begin{itemize}
\tightlist
\item
  \textbf{Not yet committed:}
  \texttt{git\ checkout\ -\/-\ content/research.md} restores the last
  committed version of that file; \texttt{git\ restore\ .} restores
  everything.
\item
  \textbf{Committed and pushed:} find the commit with
  \texttt{git\ log\ -\/-oneline}, then
  \texttt{git\ revert\ \textless{}commit-id\textgreater{}} and
  \texttt{git\ push}. This adds a new commit that undoes the old one;
  history is never lost.
\item
  \textbf{On github.com:} open the file → \textbf{History} → pick an
  older version → copy its text back.
\end{itemize}

\hypertarget{troubleshooting}{%
\section{8. Troubleshooting}\label{troubleshooting}}

\begin{longtable}[]{@{}
  >{\raggedright\arraybackslash}p{(\columnwidth - 2\tabcolsep) * \real{0.5000}}
  >{\raggedright\arraybackslash}p{(\columnwidth - 2\tabcolsep) * \real{0.5000}}@{}}
\toprule\noalign{}
\begin{minipage}[b]{\linewidth}\raggedright
Symptom
\end{minipage} & \begin{minipage}[b]{\linewidth}\raggedright
Cause and fix
\end{minipage} \\
\midrule\noalign{}
\endhead
\bottomrule\noalign{}
\endlastfoot
Actions run has a red cross & Click it, open the failed step and read
the last lines. Most often a typo in a \texttt{.toml} file (missing
quote or bracket) or Markdown header. Fix and push again. \\
\texttt{TOCSS} / SCSS error locally & Standard Hugo installed instead of
extended: \texttt{brew\ reinstall\ hugo} (Homebrew's Hugo is
extended). \\
Blank or unstyled page locally & Themes missing:
\texttt{git\ submodule\ update\ -\/-init\ -\/-recursive}. \\
\texttt{git\ push} rejected (``fetch first'') & Someone (you on
github.com) changed the repository. Run \texttt{git\ pull}, then
\texttt{git\ push}. \\
\texttt{git\ push} asks for a password & Run \texttt{gh\ auth\ login},
or use a personal access token. \\
Changes not visible online & Wait for the green tick in Actions, then
hard-reload (Cmd+Shift+R). \\
New post missing online & \texttt{draft\ =\ true} in its header; set it
to \texttt{false}. \\
Pages shows a 404 & Settings → Pages → Source must be \textbf{GitHub
Actions}. \\
\end{longtable}

\hypertarget{command-cheat-sheet}{%
\section{9. Command cheat sheet}\label{command-cheat-sheet}}

\begin{longtable}[]{@{}
  >{\raggedright\arraybackslash}p{(\columnwidth - 2\tabcolsep) * \real{0.5000}}
  >{\raggedright\arraybackslash}p{(\columnwidth - 2\tabcolsep) * \real{0.5000}}@{}}
\toprule\noalign{}
\begin{minipage}[b]{\linewidth}\raggedright
Task
\end{minipage} & \begin{minipage}[b]{\linewidth}\raggedright
Command
\end{minipage} \\
\midrule\noalign{}
\endhead
\bottomrule\noalign{}
\endlastfoot
Get latest version & \texttt{git\ pull} \\
Preview (Congo) &
\texttt{hugo\ server\ -\/-configDir\ sites/congo\ -\/-themesDir\ themes} \\
Preview (Blowfish) &
\texttt{hugo\ server\ -\/-configDir\ sites/blowfish\ -\/-themesDir\ themes} \\
Preview (PaperMod) &
\texttt{hugo\ server\ -\/-configDir\ sites/papermod\ -\/-themesDir\ themes} \\
New blog post &
\texttt{hugo\ new\ content\ posts/NAME.md\ -\/-configDir\ sites/congo\ -\/-themesDir\ themes} \\
See what changed & \texttt{git\ status} / \texttt{git\ diff} \\
Publish &
\texttt{git\ add\ .\ \&\&\ git\ commit\ -m\ "message"\ \&\&\ git\ push} \\
Update themes &
\texttt{git\ submodule\ update\ -\/-remote\ -\/-merge} \\
Undo a pushed commit &
\texttt{git\ revert\ \textless{}id\textgreater{}\ \&\&\ git\ push} \\
\end{longtable}

\hypertarget{references}{%
\section{10. References}\label{references}}

\begin{itemize}
\tightlist
\item
  Hugo documentation: \url{https://gohugo.io/documentation/}
\item
  Congo theme documentation:
  \url{https://jpanther.github.io/congo/docs/}
\item
  Blowfish theme documentation: \url{https://blowfish.page/docs/}
\item
  PaperMod wiki:
  \url{https://github.com/adityatelange/hugo-PaperMod/wiki}
\item
  GitHub Pages documentation: \url{https://docs.github.com/en/pages}
\item
  Hosting Hugo on GitHub Pages:
  \url{https://gohugo.io/host-and-deploy/host-on-github-pages/}
\item
  Markdown guide: \url{https://www.markdownguide.org/basic-syntax/}
\item
  Git handbook: \url{https://git-scm.com/book/en/v2}
\end{itemize}

\end{document}
